% ================================================================== 8
\section{Estrategia de pruebas}
\label{sec:pruebas}

\subsection{Organización de la batería de pruebas}

La suite tiene \textbf{397 pruebas automatizadas} en 15 archivos, uno por
componente. Todas pasan, con \textbf{94\,\%} de cobertura de líneas. Se ejecutan
con \texttt{python -m pytest} y en GitHub Actions (Ubuntu y Windows, Python 3.10,
3.12 y 3.13) en cada \emph{push}.

\begin{table}[H]
\centering\small
\caption{Archivos de prueba, identificadores, casos y tipo de prueba.}
\label{tab:archivos-prueba}
\begin{tabularx}{\linewidth}{@{}l l r R{2.3cm} Y@{}}
\toprule
\textbf{Archivo} & \textbf{IDs} & \textbf{Casos} & \textbf{Tipo} & \textbf{Componente}\\
\midrule
\texttt{test\_aritmetica.py} & A1--A8 & 47 & Unitarias & mcd, Bézout, inverso, exponenciaciones, ventana, contador\\
\texttt{test\_primalidad.py} & P1--P11 & 94 & Unitarias & división, Miller--Rabin, MR determinista, criba, \texttt{generar\_primo}\\
\texttt{test\_factorizacion.py} & F1--F9 & 76 & Unitarias & división, Rho, Brent, \texttt{factorizar}, ataque\\
\texttt{test\_rsa.py} & R1--R12 & 52 & Unitarias & claves manuales y FIPS 186-5, cifrado, descifrado, TCR, texto\\
\texttt{test\_formatos.py} & R9 & 19 & Unitarias & instancias, claves y resultados (JSON, CSV)\\
\texttt{test\_propiedades.py} & -- & 7 & Propiedades & Bézout, inverso, ida y vuelta RSA, factorización, primo determinista\\
\texttt{test\_integracion.py} & I1--I6 & 21 & Integración & pipeline, CLI, entradas inválidas, scripts\\
\texttt{test\_demo\_optimizacion.py} & I7 & 5 & Integración & base y optimizado dan el mismo resultado\\
\texttt{test\_experimentos.py} & X1--X6 & 12 & Integración & formato, corrección y reproducibilidad del benchmark\\
\texttt{test\_interfaz\_comun.py} & G1 & 16 & Unitarias & lectura de entradas y comparación de alternativas en la interfaz\\
\texttt{test\_interfaz\_exponenciacion.py} & G2 & 9 & Unitarias & pestaña de exponenciación y paso a paso de la binaria\\
\texttt{test\_interfaz\_primalidad.py} & G3 & 10 & Unitarias & pestaña de primalidad y generación de primos\\
\texttt{test\_interfaz\_factorizacion.py} & G4 & 10 & Unitarias & pestaña de factorización y ataque\\
\texttt{test\_interfaz\_rsa.py} & G5 & 9 & Unitarias & pestañas RSA y Optimización\\
\texttt{test\_interfaz\_ventana.py} & G6 & 10 & Integración & la ventana completa: cada botón muestra el resultado correcto\\
\midrule
\textbf{Total} & & \textbf{397} & & \\
\bottomrule
\end{tabularx}
\end{table}

\subsection{Casos normales, límite, adversos, aleatorios y de referencia}

\begin{table}[H]
\centering
\caption{Cobertura de cada componente por tipo de caso (identificadores de prueba).}
\label{tab:bateria}
\scriptsize
\renewcommand{\arraystretch}{1.25}
\begin{tabularx}{\linewidth}{@{}R{1.9cm} Y Y Y Y Y@{}}
\toprule
\textbf{Componente} & \textbf{Normal} & \textbf{Límite} & \textbf{Adverso} & \textbf{Aleatorio con semilla} & \textbf{Referencia}\\
\midrule
Aritmética & A1--A3: $17^{-1}\bmod3120=2753$ & A4: exponente 0, módulo 1; A5: $k=1..8$ & A6: negativos, módulo $\le0$, $k$ inválido & Bézout e inverso (Hypothesis) & \texttt{pow} de Python; traza manual (A2)\\
Primalidad & P5: primos conocidos & P2: $n\in\{0,1,2,3\}$; P11: límite $3{,}18\times10^{23}$ & P3: Carmichael; P4: pseudoprimos fuertes ($2047$, \dots) & P10: \texttt{generar\_primo} con semillas; Hypothesis & \textbf{exhaustiva}: $n\in[0,20\,000]$ frente a \texttt{sympy.isprime} (P1); P8: cota de $1/4$\\
Factorización & F1: $3233$, $600\,851\,475\,143$ & F6: cuadrados de primos, potencias, $2^{61}-1$ & F7: $n\le0$, método desconocido, \texttt{max\_iter} agotado & F3: enteros de 60 bits; F5: semiprimos hasta 80 bits & \textbf{exhaustiva}: $[1,3000]$ frente a \texttt{sympy.factorint} (F2); cruzada (F4)\\
RSA & R1: ejemplo manual; R3: ida y vuelta 16--2048 bits & R7: $m\in\{0,1,n-1\}$; R8: texto vacío & R10: $p=q$, $e$ no invertible, $m\ge n$, \texttt{True} & R4: claves FIPS reproducibles; Hypothesis & \textbf{exhaustiva}: los 3233 mensajes de la clave manual (R2)\\
Formatos & R9: guardar y leer & R9: solo clave pública & R9: JSON mal formado, clave alterada & -- & valores conocidos\\
Integración & I1: 6 configuraciones; I7 & I1: mensajes en el borde del rango & I4: entradas inválidas en la CLI & I1: misma clave con la misma semilla & \texttt{pow} y mensaje original\\
Interfaz & G2--G5: ejemplo del informe en cada pestaña & G2: exponente 0, módulo 1; G3: $n \ge 2^{48}$ omite la división & G1: textos no numéricos; G5: $m \ge n$, criptograma vacío & G2--G4: números aleatorios con semilla & \texttt{pow} de Python; G6: ventana completa\\
Experimentos & X1: formato de CSV & X5: carpeta sin datos & X5: experimento desconocido, configuración incompleta & X3: dos corridas, mismos conteos & X2: \texttt{sympy} y \texttt{pow} en cada fila\\
\bottomrule
\end{tabularx}
\end{table}

\lstinputlisting[style=datos, caption={Salida de \texttt{python -m pytest --cov=rsa\_edu}.}]{salidas/pytest.txt}

\subsection{Cómo se verifica que las salidas son correctas}

Se hacen dos comparaciones con propósitos distintos:

\begin{table}[H]
\centering\small
\caption{Qué se compara con qué.}
\label{tab:comparaciones}
\renewcommand{\arraystretch}{1.25}
\begin{tabularx}{\linewidth}{@{}R{2.5cm} Y Y@{}}
\toprule
& \textbf{1. Corrección} & \textbf{2. Eficiencia (optimización)}\\
\midrule
Pregunta & ¿Nuestros algoritmos dan el resultado correcto? & ¿La segunda alternativa es más rápida que la primera?\\
Se compara & Nuestro resultado contra una \textbf{implementación de referencia externa} y contra la definición & Nuestra primera alternativa contra nuestra segunda alternativa, sobre las mismas instancias\\
Referencias & \texttt{pow(b, e, m)} de Python (en C), \texttt{sympy.isprime}, \texttt{sympy.factorint}, \texttt{gmpy2.powmod}; valores conocidos (ejemplo manual, Carmichael, pseudoprimos fuertes) & Contador de operaciones y tiempo medido\\
Dónde & \texttt{tests/}, columna \texttt{ok} de cada experimento, demos & \texttt{benchmark.py}, \texttt{comparar}, \texttt{demo\_optimizacion.py}\\
\bottomrule
\end{tabularx}
\end{table}

Los algoritmos de la biblioteca no usan estas referencias: \texttt{sympy} y
\texttt{pow} con tres argumentos aparecen solo en pruebas, experimentos y demos.
La única biblioteca externa que el sistema puede usar en su ejecución es
\texttt{gmpy2}, de forma opcional en \texttt{pipeline.py}; la prueba I1 verifica
que esa configuración produce la misma clave y los mismos cifrados que la
nuestra. La eficiencia no se juzga contra una biblioteca escrita en C, porque la
diferencia reflejaría el lenguaje y no el algoritmo; por eso las dos alternativas
comparadas son implementaciones propias en las mismas condiciones.

A estas comparaciones se suman otros tres mecanismos:
\begin{itemize}
  \item \textbf{Invariantes de salida}, que no necesitan oráculo (Tabla~\ref{tab:invariantes}).
  \item \textbf{Pruebas basadas en propiedades (Hypothesis)}: se declara una
        propiedad («para todo $m$: descifrar(cifrar($m$)) $= m$») y la herramienta
        genera cientos de entradas y reduce cualquier fallo al contraejemplo mínimo.
  \item \textbf{Verificación estadística} de los algoritmos probabilísticos: la
        tasa de error de Miller--Rabin con una ronda queda bajo la cota de $1/4$ (P8).
\end{itemize}

\begin{table}[H]
\centering\small
\caption{Invariantes que se comprueban automáticamente.}
\label{tab:invariantes}
\begin{tabularx}{\textwidth}{@{}l Y@{}}
\toprule
\textbf{Módulo} & \textbf{Invariante}\\
\midrule
Euclides extendido & $a x + b y = g$, $g \mid a$ y $g \mid b$\\
Inverso modular & $a \cdot a^{-1} \equiv 1 \pmod m$ y $0 < a^{-1} < m$\\
Exponenciación & coincide con \texttt{pow}; la binaria hace a lo sumo $2\lfloor\log_2 e\rfloor + 1$ multiplicaciones; la ventana, menos\\
Factorización & $\prod F = n$, todo $f \in F$ es primo y $F$ está ordenada\\
Pollard Rho y Brent & si devuelve $d$, entonces $1 < d < n$ y $d \mid n$\\
Claves & $n = pq$, $p \ne q$ primos, $e d \equiv 1 \pmod{\lambda(n)}$; las aleatorias cumplen FIPS 186-5\\
Cifrado & $\textsc{Descifrar}(\textsc{Cifrar}(m)) = \textsc{DescifrarTCR}(\textsc{Cifrar}(m)) = m$\\
\bottomrule
\end{tabularx}
\end{table}

\subsection{Validación con la instancia resuelta a mano}

La instancia de la Sección~\ref{sec:instancia} se reproduce exactamente. El
script \texttt{scripts/traza\_instancia.py} muestra la tabla de Euclides
extendido, la traza de la exponenciación y el descifrado por TCR, y compara cada
valor con \texttt{pow}:

\lstinputlisting[style=datos, caption={Salida de \texttt{python scripts/traza\_instancia.py}.}]{salidas/traza.txt}

Durante la sesión se puede proponer cualquier instancia; el script verifica por
sí mismo cada resultado, y rechaza las inválidas:

\lstinputlisting[style=datos, caption={Salida de \texttt{python scripts/traza\_instancia.py --p 61 --q 51}.}]{salidas/inval1.txt}

\subsection{Equivalencia entre alternativas}

\lstinputlisting[style=datos, caption={Las alternativas de cada etapa sobre la misma instancia (\texttt{python -m rsa\_edu comparar}).}]{salidas/comparar.txt}

\begin{table}[H]
\centering\small
\caption{Resultados de ambas alternativas sobre las mismas instancias.}
\begin{tabularx}{\linewidth}{@{}l Y Y c@{}}
\toprule
\textbf{Instancia} & \textbf{Primera alternativa} & \textbf{Segunda alternativa} & \textbf{¿Coinciden?}\\
\midrule
$65^{17}\bmod3233$ & binaria: 2790 & ventana: 2790 & sí\\
$561$, $2047$ & división: compuestos & MR y MR determinista: compuestos & sí\\
$n=3233$ & división: $[53, 61]$ & Rho y Brent: $[53, 61]$ & sí\\
$2^{32}+1$ & división: $[641, 6700417]$ & Rho y Brent: $[641, 6700417]$ & sí\\
Descifrar $2790$ & directo: 65 & TCR: 65 & sí\\
Primo de 512 bits, semilla 7 & sin criba: $12161158388\ldots$ & con criba: $12161158388\ldots$ & sí, el mismo\\
\bottomrule
\end{tabularx}
\end{table}

\begin{table}[H]
\centering\small
\caption{Diferencias observadas entre alternativas y su explicación.}
\label{tab:diferencias}
\renewcommand{\arraystretch}{1.2}
\begin{tabularx}{\linewidth}{@{}R{4.3cm} Y R{2.6cm}@{}}
\toprule
\textbf{Diferencia} & \textbf{Explicación} & \textbf{Tipo}\\
\midrule
Binaria usa 6 multiplicaciones y ventana 5 ($e=17$) & El valor es el mismo; cambia el trabajo, que es lo que se mide. & Propiedad esperada\\
Pollard Rho devuelve 53 o 61 según la semilla & Es aleatorizado; \texttt{factorizar} ordena la lista y el resultado final es idéntico. & Decisión de implementación\\
MR con la sola base 2 declara primo a $2047=23\cdot89$ & $2047$ es pseudoprimo fuerte en base 2; con bases aleatorias o la base 3 se detecta. & Propiedad esperada\\
$d=2753$ con $\varphi(n)$; $d=413$ con $\lambda(n)$ & $\lambda(n)$ divide a $\varphi(n)$; ambos descifran $2790\rightarrow65$. & Decisión (FIPS 186-5)\\
Con $n=3233$ el TCR no es más rápido & Con claves diminutas pesa el costo fijo de reducir y recombinar; la ventaja aparece desde 512 bits. & Propiedad esperada\\
Con $\le24$ bits la división de prueba gana a MR & $\sqrt n$ es pequeño y MR tiene un costo fijo de 20 rondas. & Propiedad esperada\\
\bottomrule
\end{tabularx}
\end{table}

Ninguna diferencia corresponde a un error pendiente. La equivalencia se comprueba
automáticamente en cada \emph{push} (pruebas P1, F2, F4, A4, R2, R6, I1 e I7).

\subsection{Semillas y reproducibilidad}

Todo lo aleatorio del proyecto usa una semilla fija (2026 por defecto): las
instancias (\texttt{scripts/\allowbreak generar\_instancias.py}, que registra la semilla en
el campo \texttt{meta} de cada archivo JSON), las configuraciones de los
experimentos (\texttt{config/*.json}), las pruebas (\texttt{tests/conftest.py}) y
las funciones de la biblioteca, que reciben el generador como parámetro
(\texttt{rng}). Con la misma semilla se obtienen las mismas instancias, las
mismas claves y los mismos conteos en cualquier computadora: regenerar las
instancias produce archivos idénticos byte a byte, y las pruebas R4, X3 y la
propiedad \texttt{test\_generar\_\allowbreak primo\_\allowbreak determinista} lo verifican.

